\documentclass[11pt,letterpaper]{article}
\usepackage[T1]{fontenc}
\usepackage[utf8]{inputenc}
\usepackage{lmodern}
\usepackage[margin=0.86in]{geometry}
\usepackage{microtype}
\usepackage{booktabs,longtable,tabularx,array}
\usepackage{enumitem}
\usepackage{xcolor}
\usepackage{hyperref}
\usepackage{xurl}
\definecolor{DocBlue}{HTML}{16435F}
\hypersetup{colorlinks=true,urlcolor=DocBlue,linkcolor=DocBlue}
\setlength{\parindent}{0pt}
\setlength{\parskip}{0.54em}
\setlist[itemize]{leftmargin=*,itemsep=3pt}
\begin{document}
\begin{center}
{\LARGE\bfseries Delegated Dismissal and Exception Lifecycle}\par
\vspace{5pt}{\large Keep review authority separate from direct state mutation}\par
\vspace{9pt}{\small Engineering reference | Proposed design | Not a deployment or approval record}\par
\end{center}
\tableofcontents
\clearpage
\section{Distinct operations}
Direct dismissal uses PATCH /repos/\{owner\}/\{repo\}/code-scanning/alerts/\{alert\_number\} with state dismissed and a platform-supported reason; it changes alert state across branches. Delegated dismissal uses dismissal-request APIs when enabled and an authorized reviewer approves or denies. Neither action resolves an unrelated PR review or required check.

Request creation or direct dismissal behavior and API availability must be evaluated against the configured GitHub environment, current REST version and org permissions before implementation.


\section{Lifecycle}
\begin{itemize}
\item Submitted: developer supplies finding, reproduction and rationale.
\item Validated: read-only engine confirms fingerprint, full path, repository scope, ownership and independent evidence.
\item Review pending: authorized AppSec/security manager applies decision rights.
\item Approved: independent approver and policy version recorded with expiry or revalidation rules.
\item Executed: separate privileged executor revalidates current state and action boundary.
\item Reconciled: audit log correlates requested action, API result and observed state.
\item Expired/revoked/reopened: resume investigation and reassess merge eligibility.
\end{itemize}

\section{Approval safety}
Never approve just because a comment has a minimum length, the reason is in an allowlist or the term compensating control appears. Require a documented and independently verified security justification. Genuine unresolved vulnerabilities go through risk acceptance, not a false-positive shortcut.

Automatically processing a previously authorized narrowly scoped exception still requires proof that the new request matches approved scope. Requests without proof are routed to human review.


\section{Identity and scope}
\begin{itemize}
\item Approval and requester identities must differ and be verified against authority for the organization and repository.
\item Reject wildcards and broad path-only matches for automatic dismissal; bind alert, rule, fingerprint, component and source revision.
\item Require an expiration or documented event-driven revalidation condition.
\item Record objective evidence ID, secure link, decision timestamp and policy digest.
\end{itemize}

\section{Cross-branch implications}
GitHub documents that dismissing a code scanning alert dismisses it on all branches. Do not claim per-PR dismissal scope. Model changes to future scans and reopening behavior explicitly and validate in pilot.


\section{Governance}
Keep immutable evidence of who requested, who independently approved, why, for how long, and the exact GitHub API response. Use retention and data-classification policies; never publish real vulnerabilities or private repository inventories in this public docs repository.


\section*{Document status and provenance}
This document is a proposed reference design. Actual GitHub feature availability, organization policies, approvals, and runtime performance must be independently verified before adoption. The architecture inventory is not proof of vulnerability exploitability or correctness of a dismissal.
\end{document}
